% --- START OF FILE 09_persistenza.tex ---

\chapter{Persistenza dei Dati: Room e Firebase}

Un'app reale necessita di salvare dati in modo che sopravvivano alla chiusura o alla rotazione dello schermo.
In Android abbiamo diverse scelte:
\begin{itemize}
    \item \textbf{SharedPreferences / DataStore:} Per dati semplici chiave-valore (es.
    impostazioni, token di login).
    \item \textbf{File System:} Per dati non strutturati (es.
    foto, PDF).
    \item \textbf{Database (SQLite / Room):} Per dati strutturati, liste, relazioni e query complesse locali.
    \item \textbf{Cloud (Firebase):} Per dati condivisi tra più utenti/dispositivi.
\end{itemize}

\section{SQLite "Puro" (Il Motore)}
Android include nativamente \textbf{SQLite}, un database relazionale embedded (non c'è un processo server separato, ma una libreria C molto compatta e veloce).
Caratteristica peculiare di SQLite è il \textit{Type Affinity}: il controllo dei tipi è "rilassato" (puoi inserire una stringa in una colonna numerica).

L'approccio storico (ora considerato \textit{obsoleto} per la troppa verbosità) richiedeva:
\begin{enumerate}
    \item Estendere \texttt{SQLiteOpenHelper} per gestire creazione (\texttt{onCreate}) e aggiornamento schema (\texttt{onUpgrade}).
    \item Scrivere query SQL manuali, proni a errori a runtime (non validati dal compilatore).
    \item Usare \texttt{ContentValues} per gli \texttt{INSERT} e i \texttt{Cursor} (puntatori ai risultati) per leggere i dati estratti con query \texttt{SELECT}.
\end{enumerate}

\section{Jetpack Room: L'Approccio Moderno}
Per risolvere i problemi di SQLite puro (eccesso di codice \textit{boilerplate} e zero type-safety), Google ha introdotto \textbf{Room}, un layer di astrazione sopra SQLite.

Room è una \textbf{libreria che funziona come astrazione sopra SQLite}.
Si basa su 3 elementi fondamentali:
\begin{enumerate}
    \item \textbf{Entity:} Le data class Kotlin annotate con \texttt{@Entity}.
    In Room, una classe annotata con \texttt{@Entity} rappresenta \textbf{una tabella del database}.
    \item \textbf{DAO (Data Access Object):} Interfacce annotate con \texttt{@Dao}.
    In Room, il DAO è responsabile di \textbf{gestire le query}.
    I metodi annotati con \texttt{suspend} indicano che la funzione \textbf{deve essere chiamata in una coroutine}.
    Un metodo che ritorna un \texttt{Flow} serve a \textbf{emettere dati nel tempo} (aggiornando automaticamente la UI).
    \item \textbf{Database:} Una classe astratta annotata con \texttt{@Database}.
    Room controlla le query SQL \textbf{a compile-time}, garantendo sicurezza prima dell'esecuzione.
\end{enumerate}

\begin{center}
    \begin{tikzpicture}[
        node distance=0.8cm,
        box/.style={rectangle, draw, rounded corners, fill=white, text width=5cm, align=center, drop shadow},
        arrow/.style={-Stealth, thick}
    ]
        \node[box, fill=purple!10] (ui) {\textbf{UI (Jetpack Compose)}\\Osserva lo stato};
        \node[box, fill=blue!10, below=of ui] (vm) {\textbf{ViewModel}\\Raccoglie il Flow};
        \node[box, fill=green!10, below=of vm] (repo) {\textbf{Repository}\\Astrae la sorgente dati};
        \node[box, fill=orange!10, below=of repo] (dao) {\textbf{Room (DAO)}\\Converte oggetti in SQL};
        \node[box, fill=gray!20, below=of dao] (sqlite) {\textbf{SQLite}\\Motore fisico su disco};

        \draw[arrow] (ui) -- (vm);
        \draw[arrow] (vm) -- (repo);
        \draw[arrow] (repo) -- (dao);
        \draw[arrow] (dao) -- (sqlite);
    \end{tikzpicture}
\end{center}

\subsection{Configurazione e KSP}
Per far funzionare Room, Android deve "scrivere il codice SQLite al posto tuo" durante la compilazione.
Questo si fa usando il \textbf{KSP (Kotlin Symbol Processing)}.
Se nel file \texttt{build.gradle.kts} dimentichi di aggiungere il plugin KSP o la dipendenza \texttt{ksp(libs.androidx.room.compiler)}, il progetto non compilerà.

\subsection{Esempio di Codice Room}
\begin{lstlisting}[caption={I tre componenti di Room}]
// 1. ENTITY (Tabella)
@Entity(tableName = "shopping_items")
data class ShoppingItem(
    @PrimaryKey(autoGenerate = true) val id: Int = 0,
    val name: String,
    val quantity: Int
)

// 2. DAO (Query)
@Dao
interface ShoppingDao {
    // Flow trasforma la query in uno stream Reattivo!
    @Query("SELECT * FROM shopping_items")
    fun getAllItems(): Flow<List<ShoppingItem>>

    // suspend impone l'uso di Coroutine (no blocco UI)
    @Insert
    suspend fun insertItem(item: ShoppingItem)
}

// 3. DATABASE
@Database(entities = [ShoppingItem::class], version = 1)
abstract class AppDatabase : RoomDatabase() {
    abstract fun shoppingDao(): ShoppingDao
}
\end{lstlisting}

\begin{tcolorbox}[colback=blue!5!white,colframe=blue!75!black,title=Da Flow a UI Reattiva in Compose]
    Una potenza di Room + Compose è l'uso del \textbf{\texttt{Flow}}.
    Quando il DB viene aggiornato (es.
    inserimento), Room emette un nuovo \texttt{Flow}.
    Nella UI (Compose), usando la funzione \textbf{\texttt{collectAsStateWithLifecycle()}}, il \texttt{Flow} si converte in \texttt{State}.
    Al mutare dello State, Compose fa automaticamente la \textbf{Ricomposizione}, ridisegnando la lista senza bisogno di \texttt{notifyDataSetChanged()}!
\end{tcolorbox}

\section{Backend as a Service: Firebase}
Firebase è una piattaforma di Google che fornisce servizi di backend pronti all'uso.
Sostituisce la necessità di creare un proprio server, gestire database remoti e sviluppare API. Servizi principali: \textit{Authentication}, \textit{Cloud Storage} (per file), e i due database NoSQL.

\subsection{Realtime Database vs Cloud Firestore}
Una tematica classica da esame è la differenza tra i due database di Firebase.
Entrambi offrono sincronizzazione in tempo reale tramite \textbf{WebSocket}.

\begin{table}[h]
    \centering
    \begin{tabularx}{\textwidth}{|X|X|}
\hline
\rowcolor{gray!20} \textbf{Realtime Database} & \textbf{Cloud Firestore} \\ \hline
I dati sono memorizzati come un grande \textbf{JSON tree}. & I dati sono organizzati in \textbf{Collection} che contengono \textbf{Documenti}. Un document a sua volta può contenere \textbf{campi}. \\ \hline
Meno scalabile. Scaricare un nodo scarica \textit{tutti i dati} in esso contenuti (spesso si scaricano dati non necessari). & \textbf{Più scalabile}. Le query permettono il download mirato. Una query con più condizioni può richiedere la creazione di un \textbf{indice}. \\ \hline
Si usano listener come \texttt{ValueEventListener} per il tempo reale. & Il metodo \texttt{addSnapshotListener()} permette di ricevere aggiornamenti \textbf{in tempo reale}. \\ \hline
    \end{tabularx}
\end{table}

\subsection{Denormalizzazione (Best Practice per DB NoSQL)}
Nei database relazionali (SQL) si normalizzano i dati (es.
Join tables).
Nei database NoSQL come Firebase, la regola d'oro (per evitare di scaricare MegaBytes di messaggi solo per leggere il titolo
di una chat) è \textbf{appiattire i dati (Flattening) e Denormalizzare}.

Si evitano alberi profondi creando nodi separati al livello principale (top level nodes):
\begin{lstlisting}[language=SQL, caption={Antipattern vs Denormalizzazione}]
// SBAGLIATO (Nidificato)
"classes": {
    "class1": {
        "students": { "student1": true } // Se chiamo 'class1' scarico tutti gli studenti!
    }
}

// CORRETTO (Appiattito e Denormalizzato)
"classes": { "class1": {...dati classe} }
"students": { "student1": {...dati studente} }
"class_enrolments": {
    "class1": { "student1": true } // Tabella pivot usata solo per le query
}
\end{lstlisting}

\subsection{Lettura Dati e Listener}
Nei quiz è richiesto di riconoscere i listener appropriati:
\begin{itemize}
    \item \textbf{Realtime Database:} Si usa \texttt{ValueEventListener} (che ha i metodi \texttt{onDataChange} e \texttt{onCancelled}).
    \item \textbf{Firestore:} Per leggere una tantum si usa \texttt{.get().addOnSuccessListener}.
    Per ricevere aggiornamenti \textit{in tempo reale}, si usa \texttt{addSnapshotListener()}.
\end{itemize}

Le query in Firebase sono \textbf{Asincrone}.
Non bloccano la UI. I dati ottenuti vanno poi mappati in liste.
Se si usa il vecchio sistema XML con \texttt{SimpleAdapter}, bisogna ricordare che \textbf{non supporta nativamente il caricamento di immagini da URL}; per
farlo occorre definire un \texttt{ViewBinder} personalizzato e usare librerie come \textbf{Glide} o \textbf{Coil}.